% Docrines and decent

\chapter{Logical Semantics and Coequational Approximation}
\label{ch:logical-semantics-coequational-approximation}

\section{Purpose and Imported Logical Structure}
\label{sec:logical-semantics-standing-framework}

The logical constructions of the preceding chapters produce predicates on
execution objects.  These predicates may depend on more data than residual
behaviour alone: a downstream state, a response position, a cubical execution
cell, a finite path, or a hidden feedback state.  The purpose of this
appendix is to explain how the logical predicates already constructed in the
previous chapters determine ordinary local coequational theories of residual
behaviour.

This appendix adds no new ambient logical structure.  Each construction is
read in the standing ambient context of the chapter from which it is imported.
In particular, the coequational closure and interior constructions are read
in the ambient context of
Chapter~\ref{ch:coequational-subobjects-birkhoff}, while box modalities,
fixed-point modalities, guarded-profile modalities, cubical modalities, and
feedback-trace modalities are used only in the fragments where the
corresponding right adjoints, fixed points, and Beck--Chevalley laws were
already constructed.

The logical operations used below are the ones already constructed in the
body of the book:

\begin{enumerate}
  \item local coequational theories and derivative closure are those of
  Chapter~\ref{ch:coequational-subobjects-birkhoff};

  \item syntactic transition--output images are those of
  Chapter~\ref{ch:syntactic-transition-output-eilenberg};

  \item dynamic span diamonds, boxes, tests, and finite choice are those of
  Chapter~\ref{ch:relative-mealy-program-categories};

  \item guarded response-profile modalities are those of
  Chapter~\ref{ch:polynomial-relative-mealy-actions};

  \item cubical execution predicates and dynamic cubical modalities are those
  of Chapter~\ref{ch:cubical-mealy-execution-concurrency-types};

  \item finite-path and feedback-trace modalities are those of
  Chapter~\ref{ch:feedback-trace}.
\end{enumerate}

Thus every occurrence of
\[
  \exists_f,
  \qquad
  \forall_f,
  \qquad
  \mu,
  \qquad
  \nu
\]
below refers to a previously constructed logical operation in the relevant
chapter.  The appendix only records how these already-defined predicates are
projected to residual behaviour and then closed or interiorized to obtain
local coequational theories.

For internal categories \(A\) and \(B\), write
\[
  \mathrm{Beh}^{\mathrm{can}}_{A,B}:A\rightsquigarrow B
\]
for the canonical residual behaviour Mealy morphism, and write
\[
  \mathrm{Beh}_{A,B,0}
\]
for its carrier.  When the interfaces \(A,B\) are fixed, abbreviate this
object to
\[
  \mathrm{Beh}_0.
\]
It is equipped with anchor maps
\[
  p_{\mathrm{Beh}}:\mathrm{Beh}_0\to A_0,
  \qquad
  q_{\mathrm{Beh}}:\mathrm{Beh}_0\to B_0.
\]
Set
\[
  Z:=A_0\times B_0.
\]

A \textbf{behaviour predicate} is a subobject of
\(\mathrm{Beh}_0\) in the slice
\[
  \mathcal E/Z.
\]
A \textbf{local coequational theory} is a derivative-closed behaviour
predicate, equivalently a strict sub-Mealy object of
\[
  \mathrm{Beh}^{\mathrm{can}}_{A,B}.
\]

If
\[
  \Phi:P:A\rightsquigarrow B
\]
is a Mealy morphism and
\[
  r_Y:Y\to B_0
\]
is a right \(B\)-action, the relative execution state object is
\[
  S_{\Phi,Y}:=P\times_{q_P,B_0,r_Y}Y.
\]
Its canonical instance is
\[
  S^{\mathrm{can}}_Y
  :=
  \mathrm{Beh}_0\times_{q_{\mathrm{Beh}},B_0,r_Y}Y.
\]
Residual semantics of \(P\) gives the strict semantic homomorphism
\[
  \beta_P:P\to \mathrm{Beh}^{\mathrm{can}}_{A,B},
\]
with carrier map
\[
  \beta_P:P\to \mathrm{Beh}_0.
\]
It induces a comparison map
\[
  \beta_{P,Y}:S_{\Phi,Y}\longrightarrow S^{\mathrm{can}}_Y.
\]

The comparison between canonical execution predicates and behaviour
predicates is carried by the projection
\[
  \pi_Y:S^{\mathrm{can}}_Y\longrightarrow \mathrm{Beh}_0.
\]
For a canonical execution predicate
\[
  \psi\hookrightarrow S^{\mathrm{can}}_Y,
\]
the existential and universal behaviour projections are the predicate
projections already available in the logical fragments where they occur:
\[
  \exists_{\pi_Y}\psi\hookrightarrow \mathrm{Beh}_0,
  \qquad
  \forall_{\pi_Y}\psi\hookrightarrow \mathrm{Beh}_0.
\]
The first says that some downstream state over the behaviour satisfies
\(\psi\).  The second says that every downstream state over the behaviour
satisfies \(\psi\).

These projected predicates need not be derivative-closed.  Derivative closure
and derivative interior turn them into local coequational theories:
\[
  \mathrm{May}_Y(\psi)
  :=
  \mathrm{Der}_0(\exists_{\pi_Y}\psi),
\]
and
\[
  \mathrm{Must}_Y(\psi)
  :=
  \mathrm{DerInt}(\forall_{\pi_Y}\psi).
\]

The central comparison is the adjoint triple
\[
  \mathrm{May}_Y
  \dashv
  J_Y
  \dashv
  \mathrm{Must}_Y,
  \qquad
  J_Y(V)=\pi_Y^*\,iV,
\]
where
\[
  i:\mathbf{Coeq}_{A,B}
  \hookrightarrow
  \operatorname{Sub}_{\mathcal E/Z}(\mathrm{Beh}_0)
\]
is the inclusion of local coequational theories into behaviour predicates.
The rest of the appendix shows how the execution logics constructed earlier
feed into this coequational comparison.

For an object \(X\to Z\), write
\[
  \operatorname{Pred}_Z(X)
  :=
  \operatorname{Sub}_{\mathcal E/Z}(X).
\]
When no confusion is possible, write simply
\[
  \operatorname{Pred}(X).
\]

\section{Recall: Predicate Transformers Used Earlier}
\label{sec:base-subobject-calculus}

This appendix uses the subobject semantics developed in the earlier logical
chapters.  We recall only the clauses needed for the synthesis.

A formula in context denotes a subobject of its context object.  Substitution
along a map is pullback.  Dynamic diamonds are interpreted by existential
image along the source map of a span.  Dynamic boxes are interpreted by
universal image along the same source map in the box fragments constructed in
the book.

For a span
\[
  X\xleftarrow{\ell}U\xrightarrow{r}Y
\]
from one of the dynamic fragments, the associated predicate transformers are
\[
  \langle U\rangle:
  \operatorname{Pred}(Y)\to\operatorname{Pred}(X),
  \qquad
  \langle U\rangle\varphi
  :=
  \exists_{\ell}r^*\varphi,
\]
and, in the box fragment,
\[
  [U]:
  \operatorname{Pred}(Y)\to\operatorname{Pred}(X),
  \qquad
  [U]\varphi
  :=
  \forall_{\ell}r^*\varphi.
\]

The substitution and functoriality laws used below are the Beck--Chevalley
and adjunction laws established for the corresponding earlier fragments.
No additional logical rules are added here.

\begin{proposition}[Soundness inherited from the earlier fragments]
\label{prop:base-subobject-calculus-soundness}
Every derivable entailment in any of the logical fragments recalled above is
interpreted by an inclusion of subobjects in the appropriate execution
context.
\end{proposition}

\begin{proof}
\leavevmode
\begin{enumerate}
  \item \textbf{Reduce to the imported soundness theorem.}
  Each fragment used in this appendix has already been given a subobject
  semantics in the corresponding earlier chapter.  The proposition is the
  soundness theorem for whichever fragment supplies the formula under
  consideration.

  \item \textbf{Interpret structural rules.}
  The propositional rules are interpreted by the order-theoretic operations
  on subobjects, and substitution is interpreted by pullback.

  \item \textbf{Interpret modal rules.}
  Diamonds and boxes are interpreted by the adjunctions
  \[
    \exists_{\ell}\dashv \ell^*,
    \qquad
    \ell^*\dashv \forall_{\ell}
  \]
  defining existential and universal projection along the source maps of the
  spans used in the fragment.

  \item \textbf{Use Beck--Chevalley for substitution.}
  The Beck--Chevalley laws established in the earlier fragments give the
  substitution laws for the modal operations.  Hence derivable entailments are
  interpreted by inclusions of subobjects.
\end{enumerate}
\end{proof}

\section{Canonical Relative Execution}
\label{sec:canonical-relative-execution}

Relative execution places a Mealy morphism in a downstream environment.  A
state of the resulting execution object consists of a system state and a
downstream state whose anchors agree.  During one primitive execution step,
the Mealy morphism consumes an input arrow, updates its internal state, and
produces an output-comparison arrow in \(B\).  The right \(B\)-action then
uses that output-comparison arrow to update the downstream state.

Let
\[
  \Phi:P:A\rightsquigarrow B
\]
be a Mealy morphism.  Write
\[
  p_P:P\to A_0,
  \qquad
  q_P:P\to B_0
\]
for its input and output anchors.  The transition and output maps are
\[
  \delta_P:
  A_1\times_{t_A,A_0,p_P}P\longrightarrow P,
\]
and
\[
  \omega_P:
  A_1\times_{t_A,A_0,p_P}P\longrightarrow B_1,
\]
with typing
\[
  t_B\omega_P(a,x)=q_P(x),
  \qquad
  s_B\omega_P(a,x)=q_P(\delta_P(a,x)).
\]
Thus an input arrow
\[
  a:u\to v
\]
acts on states over \(v\), and the output comparison
\[
  \omega_P(a,x)
\]
points from the new output anchor to the old output anchor.

Let
\[
  r_Y:Y\to B_0
\]
be a right \(B\)-action.  Its action map is
\[
  \sigma_Y:
  B_1\times_{t_B,B_0,r_Y}Y\longrightarrow Y,
\]
with
\[
  r_Y\sigma_Y=s_B\pi_{B_1}.
\]
Therefore a \(B\)-arrow
\[
  b:b'\to b
\]
acts on a state over \(b\) and produces a state over \(b'\).

\begin{definition}[Relative execution state object]
\label{def:relative-execution-state-object}
The \textbf{relative execution state object} of \(\Phi\) against \(Y\) is
\[
  S_{\Phi,Y}
  :=
  P\times_{q_P,B_0,r_Y}Y.
\]
Its \textbf{canonical relative execution state object} is
\[
  S^{\mathrm{can}}_Y
  :=
  \mathrm{Beh}_0\times_{q_{\mathrm{Beh}},B_0,r_Y}Y.
\]
\end{definition}

The primitive execution apex of \((\Phi,Y)\) is
\[
  M_{\Phi,Y}
  :=
  A_1\times_{t_A,A_0,p_P\pi_P}S_{\Phi,Y}.
\]
The primitive execution span is
\[
  m_{\Phi,Y}
  =
  \left(
  S_{\Phi,Y}
  \xleftarrow{s_{\Phi,Y}}
  M_{\Phi,Y}
  \xrightarrow{t_{\Phi,Y}}
  S_{\Phi,Y}
  \right),
\]
where
\[
  s_{\Phi,Y}(a,x,y)=(x,y),
\]
and
\[
  t_{\Phi,Y}(a,x,y)
  =
  \left(
    \delta_P(a,x),
    \sigma_Y(\omega_P(a,x),y)
  \right).
\]
This target map is well-typed because
\[
  t_B\omega_P(a,x)=q_P(x)=r_Y(y)
\]
and hence
\[
  r_Y\sigma_Y(\omega_P(a,x),y)
  =
  s_B\omega_P(a,x)
  =
  q_P(\delta_P(a,x)).
\]

The canonical primitive execution span is
\[
  m_{\mathrm{Beh},Y}
  =
  \left(
  S^{\mathrm{can}}_Y
  \xleftarrow{s_{\mathrm{Beh},Y}}
  M_{\mathrm{Beh},Y}
  \xrightarrow{t_{\mathrm{Beh},Y}}
  S^{\mathrm{can}}_Y
  \right),
\]
where
\[
  M_{\mathrm{Beh},Y}
  :=
  A_1\times_{t_A,A_0,p_{\mathrm{Beh}}\pi_{\mathrm{Beh}}}
  S^{\mathrm{can}}_Y.
\]

Residual semantics replaces a concrete state \(x\in P\) by its residual
behaviour \(\beta_P(x)\).  Since the downstream state \(y\) is already typed
over the same output anchor, this replacement extends fibrewise to relative
execution states:
\[
  \beta_{P,Y}
  :=
  \beta_P\times_{B_0}1_Y:
  S_{\Phi,Y}
  \longrightarrow
  S^{\mathrm{can}}_Y.
\]

The primitive residual comparison is summarized by the diagram
\[
\begin{tikzcd}[column sep=large,row sep=large]
A_1\times_{A_0}S_{\Phi,Y}
  \arrow[r, "\overline{\beta}_{P,Y}"]
  \arrow[d, shift right=.65ex, "s_{\Phi,Y}"']
  \arrow[d, shift left=.65ex, "t_{\Phi,Y}"]
&
A_1\times_{A_0}S^{\mathrm{can}}_Y
  \arrow[d, shift right=.65ex, "s_{\mathrm{Beh},Y}"']
  \arrow[d, shift left=.65ex, "t_{\mathrm{Beh},Y}"]
\\
S_{\Phi,Y}
  \arrow[r, "\beta_{P,Y}"']
&
S^{\mathrm{can}}_Y .
\end{tikzcd}
\]

\begin{proposition}[Primitive residual comparison]
\label{prop:primitive-residual-comparison}
The map
\[
  \beta_{P,Y}:S_{\Phi,Y}\to S^{\mathrm{can}}_Y
\]
extends to a morphism of primitive execution spans
\[
  m_{\Phi,Y}\longrightarrow m_{\mathrm{Beh},Y}.
\]
The square over source maps
\[
\begin{tikzcd}
M_{\Phi,Y}
  \arrow[r, "\overline{\beta}_{P,Y}"]
  \arrow[d, "s_{\Phi,Y}"']
&
M_{\mathrm{Beh},Y}
  \arrow[d, "s_{\mathrm{Beh},Y}"]
\\
S_{\Phi,Y}
  \arrow[r, "\beta_{P,Y}"']
&
S^{\mathrm{can}}_Y
  \arrow[ul, phantom, very near start, "\lrcorner"]
\end{tikzcd}
\]
is a pullback.
\end{proposition}

\begin{proof}
\leavevmode
\begin{enumerate}
  \item \textbf{Define the apex comparison.}
  Define
  \[
    \overline{\beta}_{P,Y}(a,x,y)
    :=
    (a,\beta_P(x),y).
  \]
  This is well-typed because \(\beta_P\) preserves output anchors:
  \[
    q_{\mathrm{Beh}}\beta_P=q_P.
  \]

  \item \textbf{Check source compatibility.}
  The source maps commute because
  \[
    s_{\mathrm{Beh},Y}(a,\beta_P(x),y)
    =
    (\beta_P(x),y)
    =
    \beta_{P,Y}(s_{\Phi,Y}(a,x,y)).
  \]

  \item \textbf{Identify the source square as a pullback.}
  Both apexes are obtained by pulling back the same object \(A_1\) along the
  input anchor of the corresponding state object:
  \[
  \begin{tikzcd}
  M_{\Phi,Y}
    \arrow[r]
    \arrow[d]
  &
  A_1
    \arrow[d, "t_A"]
  \\
  S_{\Phi,Y}
    \arrow[r, "p_P\pi_P"']
  &
  A_0
  \end{tikzcd}
  \qquad
  \begin{tikzcd}
  M_{\mathrm{Beh},Y}
    \arrow[r]
    \arrow[d]
  &
  A_1
    \arrow[d, "t_A"]
  \\
  S^{\mathrm{can}}_Y
    \arrow[r, "p_{\mathrm{Beh}}\pi_{\mathrm{Beh}}"']
  &
  A_0 .
  \end{tikzcd}
  \]
  Since \(p_{\mathrm{Beh}}\beta_P=p_P\), the comparison square over source
  maps is the induced pullback square.

  \item \textbf{Check target compatibility.}
  Residual semantics gives
  \[
    \beta_P(\delta_P(a,x))
    =
    \partial_a(\beta_P(x)),
  \]
  and output preservation gives
  \[
    \omega_P(a,x)
    =
    \omega_{\mathrm{Beh}}(a,\beta_P(x)).
  \]
  Therefore
  \[
  \begin{aligned}
    \beta_{P,Y}(t_{\Phi,Y}(a,x,y))
    &=
    \left(
      \beta_P(\delta_P(a,x)),
      \sigma_Y(\omega_P(a,x),y)
    \right)                                      \\
    &=
    \left(
      \partial_a(\beta_P(x)),
      \sigma_Y(\omega_{\mathrm{Beh}}(a,\beta_P(x)),y)
    \right)                                      \\
    &=
    t_{\mathrm{Beh},Y}(a,\beta_P(x),y).
  \end{aligned}
  \]

  \item \textbf{Conclude span compatibility.}
  The apex comparison commutes with both source and target maps.  Hence it
  defines a morphism of primitive execution spans.
\end{enumerate}
\end{proof}

\begin{definition}[Residual-pullback span scheme]
\label{def:residual-pullback-span-scheme}
A \textbf{residual-pullback span scheme} over \((A,B,Y)\) is a span
construction from one of the execution fragments of the book together with
its canonical residual comparison data.

Concretely, for every Mealy morphism
\[
  \Phi:P:A\rightsquigarrow B,
\]
it consists of a span
\[
  X_{\Phi,Y}
  \xleftarrow{\ell_{\Phi}}
  U_{\Phi}
  \xrightarrow{r_{\Phi}}
  Y_{\Phi,Y},
\]
the corresponding canonical span
\[
  X_{\mathrm{can},Y}
  \xleftarrow{\ell_{\mathrm{can}}}
  U_{\mathrm{can}}
  \xrightarrow{r_{\mathrm{can}}}
  Y_{\mathrm{can},Y},
\]
and residual comparison maps
\[
  \beta_X:X_{\Phi,Y}\to X_{\mathrm{can},Y},
  \qquad
  \beta_Y:Y_{\Phi,Y}\to Y_{\mathrm{can},Y},
\]
together with an apex comparison
\[
  \overline{\beta}_U:U_{\Phi}\to U_{\mathrm{can}}
\]
such that
\[
  \ell_{\mathrm{can}}\overline{\beta}_U
  =
  \beta_X\ell_{\Phi},
  \qquad
  r_{\mathrm{can}}\overline{\beta}_U
  =
  \beta_Y r_{\Phi},
\]
and the source square
\[
\begin{tikzcd}
U_{\Phi}
  \arrow[r, "\overline{\beta}_U"]
  \arrow[d, "\ell_{\Phi}"']
&
U_{\mathrm{can}}
  \arrow[d, "\ell_{\mathrm{can}}"]
\\
X_{\Phi,Y}
  \arrow[r, "\beta_X"']
&
X_{\mathrm{can},Y}
  \arrow[ul, phantom, very near start, "\lrcorner"]
\end{tikzcd}
\]
is a pullback.

The pullback square over source maps is the square to which the
Beck--Chevalley law for the available predicate transformer applies:
regular-image Beck--Chevalley in diamond fragments, and right-adjoint
Beck--Chevalley in the box fragments.
\end{definition}

The phrase ``residual-pullback'' records the mechanism used by the semantics.
The pullback square over source maps is the reason that modal truth in a
concrete execution context is obtained by pulling back modal truth from the
canonical execution context.

\begin{definition}[Residual-stable canonical fragment]
\label{def:residual-stable-canonical-fragment}
The \textbf{residual-stable canonical fragment} over \((A,B,Y)\) is generated
inside the already-constructed logical fragments by the following clauses.

\begin{enumerate}
  \item Atomic predicates are canonical predicates on canonical execution
  contexts.

  \item The class is closed under the propositional connectives present in
  the relevant fragment.

  \item If \(U\) is a residual-pullback span scheme and \(\varphi\) is in
  the fragment, then the diamond formula
  \[
    \langle U\rangle\varphi
  \]
  is in the fragment.

  \item If the corresponding box modality is part of the earlier fragment,
  then
  \[
    [U]\varphi
  \]
  is in the fragment.

  \item If the corresponding finite-path or feedback-trace modality is part
  of the earlier fragment, then its use is restricted to the path-star or
  trace constructions for which residual comparison was already shown to be
  compatible with the imported fixed-point semantics.
\end{enumerate}
\end{definition}

\begin{theorem}[Canonical pullback semantics]
\label{thm:canonical-pullback-semantics}
Let \(\varphi\) be a formula in the residual-stable canonical fragment.  Then
for every Mealy morphism
\[
  \Phi:P:A\rightsquigarrow B
\]
one has
\[
  \llbracket\varphi\rrbracket_{\Phi,Y}
  =
  \beta_X^*
  \llbracket\varphi\rrbracket_{\mathrm{can},Y},
\]
where \(\beta_X\) is the residual comparison map of the execution context in
which \(\varphi\) is interpreted.

For formulas involving finite-path or trace iteration, the statement is read
in the corresponding path-star fragment, where residual comparison is
source-cartesian for the underlying endospans and the fixed-point predicate
transformers commute with residual pullback.
\end{theorem}

\begin{proof}
\leavevmode
\begin{enumerate}
  \item \textbf{Prove the claim by induction on formulas.}
  The proof is by induction on the construction of \(\varphi\).

  \item \textbf{Handle atomic predicates.}
  Atomic canonical predicates are, by definition, interpreted in the
  \(\Phi\)-context by pullback from the canonical context.  Thus the claim
  holds for atomic formulas.

  \item \textbf{Handle propositional structure.}
  Pullback preserves the propositional structure available in the fragment.
  Hence the claim is preserved by the propositional connectives.

  \item \textbf{Handle diamonds.}
  For a diamond along a residual-pullback span scheme, write the canonical
  span as
  \[
    X_{\mathrm{can},Y}
    \xleftarrow{\ell_{\mathrm{can}}}
    U_{\mathrm{can}}
    \xrightarrow{r_{\mathrm{can}}}
    Y_{\mathrm{can},Y},
  \]
  and the corresponding \(\Phi\)-span as
  \[
    X_{\Phi,Y}
    \xleftarrow{\ell_{\Phi}}
    U_{\Phi}
    \xrightarrow{r_{\Phi}}
    Y_{\Phi,Y}.
  \]
  The comparison data form the diagram
  \[
  \begin{tikzcd}[column sep=large,row sep=large]
  U_{\Phi}
    \arrow[r, "\overline{\beta}_U"]
    \arrow[d, "\ell_{\Phi}"']
    \arrow[dr, phantom, very near start, "\lrcorner"]
  &
  U_{\mathrm{can}}
    \arrow[d, "\ell_{\mathrm{can}}"]
    \arrow[r, "r_{\mathrm{can}}"]
  &
  Y_{\mathrm{can},Y}
  \\
  X_{\Phi,Y}
    \arrow[r, "\beta_X"']
  &
  X_{\mathrm{can},Y}
  &
  Y_{\Phi,Y}
    \arrow[u, "\beta_Y"']
    \arrow[ul, "r_{\Phi}"']
  \end{tikzcd}
  \]
  with the source square a pullback and the target maps commuting with
  residual comparison.  Beck--Chevalley in the relevant diamond fragment
  gives
  \[
    \beta_X^*
    \exists_{\ell_{\mathrm{can}}}
    r_{\mathrm{can}}^*(\theta)
    =
    \exists_{\ell_{\Phi}}
    r_{\Phi}^*
    \beta_Y^*(\theta).
  \]
  The induction hypothesis applies to \(\theta\).

  \item \textbf{Handle boxes.}
  In a box fragment, the same source pullback square and the right-adjoint
  Beck--Chevalley law give
  \[
    \beta_X^*
    \forall_{\ell_{\mathrm{can}}}
    r_{\mathrm{can}}^*(\theta)
    =
    \forall_{\ell_{\Phi}}
    r_{\Phi}^*
    \beta_Y^*(\theta).
  \]
  The induction hypothesis again applies to \(\theta\).

  \item \textbf{Handle finite-path and trace modalities.}
  For finite-path and trace modalities, source-cartesian residual comparison
  identifies the one-step predicate transformer before and after pullback.
  Since pullback preserves the joins and meets used in the least and greatest
  fixed-point interpretations in the imported fragment, the induced
  fixed-point predicates are again pulled back from the canonical context.

  \item \textbf{Conclude the induction.}
  All formula constructors in the residual-stable canonical fragment preserve
  the asserted pullback form.  Therefore
  \[
    \llbracket\varphi\rrbracket_{\Phi,Y}
    =
    \beta_X^*
    \llbracket\varphi\rrbracket_{\mathrm{can},Y}.
  \]
\end{enumerate}
\end{proof}

\begin{corollary}[Primitive canonical pullback semantics]
\label{cor:primitive-canonical-pullback-semantics}
The primitive execution modality associated to
\[
  m_{\Phi,Y}:S_{\Phi,Y}\rightsquigarrow S_{\Phi,Y}
\]
is pulled back from the canonical primitive execution modality along
\[
  \beta_{P,Y}:S_{\Phi,Y}\to S^{\mathrm{can}}_Y.
\]
\end{corollary}

\begin{proof}
\leavevmode
\begin{enumerate}
  \item \textbf{Apply the primitive residual-pullback scheme.}
  Proposition~\ref{prop:primitive-residual-comparison} shows that the
  primitive execution span is a residual-pullback span scheme.

  \item \textbf{Invoke canonical pullback semantics.}
  The result is Theorem~\ref{thm:canonical-pullback-semantics} applied to
  that primitive span scheme.
\end{enumerate}
\end{proof}

\section{Coequational Closure and Interior}
\label{sec:coequational-closure-and-interior}

All subobjects of \(\mathrm{Beh}_0\) in this section are taken in the slice
\[
  \mathcal E/Z.
\]

A behaviour predicate may fail to be stable under future input.  A local
coequational theory is precisely a behaviour predicate that is closed under
all residual derivatives.  The operator \(\mathrm{Der}_0\) gives the least
derivative-closed predicate containing a given predicate.  The operator
\(\mathrm{DerInt}\) gives the largest derivative-closed predicate contained
in a given predicate.

A local coequational theory over \(A,B\) is a subobject
\[
  V\hookrightarrow \mathrm{Beh}_0
\]
whose canonical derivative structure restricts.  Equivalently, if
\[
  \delta_{\mathrm{Beh}}:
  A_1\times_{t_A,A_0,p_{\mathrm{Beh}}}
  \mathrm{Beh}_0
  \longrightarrow
  \mathrm{Beh}_0
\]
is the canonical derivative map, then \(V\) is local coequational precisely
when there exists a map
\[
  \delta_V:
  A_1\times_{t_A,A_0,p_V}V
  \longrightarrow
  V
\]
making the square
\[
\begin{tikzcd}[column sep=large,row sep=large]
A_1\times_{t_A,A_0,p_V}V
  \arrow[r, hook]
  \arrow[d, "\delta_V"']
&
A_1\times_{t_A,A_0,p_{\mathrm{Beh}}}\mathrm{Beh}_0
  \arrow[d, "\delta_{\mathrm{Beh}}"]
\\
V
  \arrow[r, hook]
&
\mathrm{Beh}_0
\end{tikzcd}
\]
commute.

The derivative domain is regarded over
\[
  Z=A_0\times B_0
\]
by the post-derivative state.  Thus a generalized element \((a,H)\), with
\[
  a:u\to v
\]
and \(H\) over \(v\), lies over
\[
  (u,q_{\mathrm{Beh}}(\partial_aH)).
\]
No additional closure condition is imposed on the output-comparison map: its
codomain remains \(B_1\), and derivative closure supplies exactly the typing
needed for the restricted transition-output structure.

We write
\[
  \mathbf{Coeq}_{A,B}
\]
for the poset of local coequational theories, and
\[
  i:\mathbf{Coeq}_{A,B}
  \hookrightarrow
  \operatorname{Sub}_{\mathcal E/Z}(\mathrm{Beh}_0)
\]
for the inclusion.

By the coequational closure theorem from
Chapter~\ref{ch:coequational-subobjects-birkhoff}, the inclusion \(i\) has
a left adjoint
\[
  \mathrm{Der}_0
  :
  \operatorname{Sub}_{\mathcal E/Z}(\mathrm{Beh}_0)
  \longrightarrow
  \mathbf{Coeq}_{A,B}.
\]
Thus, for every behaviour predicate \(K\) and every local coequational theory
\(V\),
\[
  \mathrm{Der}_0(K)\leq V
  \quad\Longleftrightarrow\quad
  K\leq iV.
\]

\begin{lemma}[Derivative interior]
\label{lem:derivative-interior}
The inclusion
\[
  i:\mathbf{Coeq}_{A,B}
  \hookrightarrow
  \operatorname{Sub}_{\mathcal E/Z}(\mathrm{Beh}_0)
\]
has a right adjoint
\[
  \mathrm{DerInt}:
  \operatorname{Sub}_{\mathcal E/Z}(\mathrm{Beh}_0)
  \longrightarrow
  \mathbf{Coeq}_{A,B}.
\]
It is given by
\[
  \mathrm{DerInt}(K)
  =
  \bigvee
  \left\{
    W\hookrightarrow \mathrm{Beh}_0
    \ \middle|\
    W\leq K,\ W\in\mathbf{Coeq}_{A,B}
  \right\}.
\]
\end{lemma}

\begin{proof}
\leavevmode
\begin{enumerate}
  \item \textbf{Specify the join.}
  The displayed join is taken in
  \[
    \operatorname{Sub}_{\mathcal E/Z}(\mathrm{Beh}_0).
  \]
  The indexing family is a family of subobjects of the fixed object
  \(\mathrm{Beh}_0\) in the fixed slice \(\mathcal E/Z\), hence is small in
  the ambient universe used in
  Chapter~\ref{ch:coequational-subobjects-birkhoff}.

  \item \textbf{Show the join is coequational.}
  By the join-closure theorem for local coequational theories in
  Chapter~\ref{ch:coequational-subobjects-birkhoff}, small joins of local
  coequational theories are again local coequational theories.  Hence
  \[
    \mathrm{DerInt}(K)
  \]
  is a local coequational theory.

  \item \textbf{Show containment in \(K\).}
  Since \(K\) is an upper bound for every local coequational theory
  \(W\leq K\), the join satisfies
  \[
    i\,\mathrm{DerInt}(K)\leq K.
  \]

  \item \textbf{Verify the universal property.}
  For a local coequational theory \(V\), the inequality
  \[
    iV\leq K
  \]
  holds exactly when \(V\) belongs to the indexing family of the displayed
  join.  This is equivalent to
  \[
    V\leq \mathrm{DerInt}(K).
  \]
  Therefore
  \[
    i\dashv \mathrm{DerInt}.
  \]
\end{enumerate}
\end{proof}

Consequently, for every local coequational theory \(V\) and every behaviour
predicate \(K\),
\[
  iV\leq K
  \quad\Longleftrightarrow\quad
  V\leq \mathrm{DerInt}(K).
\]

\section{May and Must Coequational Approximation}
\label{sec:may-must-coequational-approximation}

Let
\[
  Y\to B_0
\]
be a right \(B\)-action.  The canonical relative execution object
\[
  S^{\mathrm{can}}_Y
  =
  \mathrm{Beh}_0\times_{B_0}Y
\]
is an object of the slice \(\mathcal E/Z\) by projection to
\[
  \mathrm{Beh}_0\to Z.
\]
The projection
\[
  \pi_Y:S^{\mathrm{can}}_Y\to \mathrm{Beh}_0
\]
is a morphism in \(\mathcal E/Z\).

An execution predicate
\[
  \psi\hookrightarrow S^{\mathrm{can}}_Y
\]
may depend both on residual behaviour and on downstream state.  To compare it
with a local coequational theory, the downstream state must be removed.  The
earlier logical semantics supplies two projection operations along
\(\pi_Y\).  Existential projection keeps a residual behaviour when some
compatible downstream state satisfies \(\psi\).  Universal projection keeps a
residual behaviour when every compatible downstream state satisfies \(\psi\).
Derivative closure and derivative interior then force these projected
behaviour predicates to be stable under residual derivatives.

\begin{definition}[Existential and universal behaviour projections]
\label{def:existential-universal-behaviour-projections}
For a canonical execution predicate
\[
  \psi\hookrightarrow S^{\mathrm{can}}_Y
\]
in \(\mathcal E/Z\), define its \textbf{existential behaviour projection} by
\[
  \Diamond_Y\psi
  :=
  \exists_{\pi_Y}\psi
  \hookrightarrow
  \mathrm{Beh}_0.
\]
Define its \textbf{universal behaviour projection} by
\[
  \Box_Y\psi
  :=
  \forall_{\pi_Y}\psi
  \hookrightarrow
  \mathrm{Beh}_0.
\]
\end{definition}

\begin{definition}[May- and must-coequationalization]
\label{def:may-must-coequationalization}
For a canonical execution predicate
\[
  \psi\hookrightarrow S^{\mathrm{can}}_Y,
\]
define its \textbf{may-coequationalization} by
\[
  \mathrm{May}_Y(\psi)
  :=
  \mathrm{Der}_0(\exists_{\pi_Y}\psi).
\]
Define its \textbf{must-coequationalization} by
\[
  \mathrm{Must}_Y(\psi)
  :=
  \mathrm{DerInt}(\forall_{\pi_Y}\psi).
\]
\end{definition}

\begin{definition}[Execution extension of a coequational theory]
\label{def:execution-extension-coequational-theory}
For a local coequational theory
\[
  V\hookrightarrow \mathrm{Beh}_0,
\]
define its \textbf{canonical \(Y\)-extension} by
\[
  J_Y(V)
  :=
  \pi_Y^*iV
  =
  V\times_{B_0}Y
  \hookrightarrow
  S^{\mathrm{can}}_Y.
\]
This gives a monotone map
\[
  J_Y:\mathbf{Coeq}_{A,B}
  \longrightarrow
  \operatorname{Pred}_Z(S^{\mathrm{can}}_Y).
\]
\end{definition}

The three comparison operations fit into the diagram
\[
\begin{tikzcd}[column sep=large,row sep=large]
\operatorname{Pred}_Z(S^{\mathrm{can}}_Y)
  \arrow[r, shift left=1.1ex, "\exists_{\pi_Y}"]
  \arrow[r, shift right=1.1ex, "\forall_{\pi_Y}"']
&
\operatorname{Pred}_Z(\mathrm{Beh}_0)
  \arrow[r, shift left=1.1ex, "\mathrm{Der}_0"]
  \arrow[r, shift right=1.1ex, "\mathrm{DerInt}"']
&
\mathbf{Coeq}_{A,B}
\\
\mathbf{Coeq}_{A,B}
  \arrow[u, "J_Y"]
  \arrow[urr, phantom, "\dashv"]
\end{tikzcd}
\]
where the displayed arrows encode the two composites defining
\(\mathrm{May}_Y\) and \(\mathrm{Must}_Y\).

\begin{theorem}[May--must coequational approximation]
\label{thm:may-must-coequational-approximation}
There is an adjoint triple of posets
\[
  \mathrm{May}_Y
  \dashv
  J_Y
  \dashv
  \mathrm{Must}_Y.
\]
Equivalently, for every canonical execution predicate
\[
  \psi\hookrightarrow S^{\mathrm{can}}_Y
\]
and every local coequational theory
\[
  V\hookrightarrow \mathrm{Beh}_0,
\]
one has
\[
  \mathrm{May}_Y(\psi)\leq V
  \quad\Longleftrightarrow\quad
  \psi\leq J_Y(V),
\]
and
\[
  J_Y(V)\leq\psi
  \quad\Longleftrightarrow\quad
  V\leq \mathrm{Must}_Y(\psi).
\]
\end{theorem}

\begin{proof}
\leavevmode
\begin{enumerate}
  \item \textbf{Prove the left adjunction.}
  The inequality
  \[
    \mathrm{May}_Y(\psi)\leq V
  \]
  is equivalent to
  \[
    \mathrm{Der}_0(\exists_{\pi_Y}\psi)\leq V.
  \]
  By
  \[
    \mathrm{Der}_0\dashv i,
  \]
  this is equivalent to
  \[
    \exists_{\pi_Y}\psi\leq iV.
  \]
  By
  \[
    \exists_{\pi_Y}\dashv \pi_Y^*,
  \]
  this is equivalent to
  \[
    \psi\leq \pi_Y^*iV=J_Y(V).
  \]
  Therefore
  \[
    \mathrm{May}_Y\dashv J_Y.
  \]

  \item \textbf{Prove the right adjunction.}
  The inequality
  \[
    J_Y(V)\leq\psi
  \]
  is
  \[
    \pi_Y^*iV\leq\psi.
  \]
  By
  \[
    \pi_Y^*\dashv\forall_{\pi_Y},
  \]
  this is equivalent to
  \[
    iV\leq\forall_{\pi_Y}\psi.
  \]
  By
  \[
    i\dashv\mathrm{DerInt},
  \]
  this is equivalent to
  \[
    V\leq \mathrm{DerInt}(\forall_{\pi_Y}\psi)
    =
    \mathrm{Must}_Y(\psi).
  \]
  Therefore
  \[
    J_Y\dashv\mathrm{Must}_Y.
  \]

  \item \textbf{Assemble the triple.}
  Combining the two adjunctions gives
  \[
    \mathrm{May}_Y
    \dashv
    J_Y
    \dashv
    \mathrm{Must}_Y.
  \]
\end{enumerate}
\end{proof}

\begin{remark}[Vacuity of universal projections]
\label{rem:vacuity-universal-projections}
The predicate
\[
  \forall_{\pi_Y}\psi
\]
is fibrewise universal along
\[
  \pi_Y:S^{\mathrm{can}}_Y\to\mathrm{Beh}_0.
\]
Thus a behaviour with empty \(\pi_Y\)-fibre satisfies it vacuously.  Write
\[
  \operatorname{supp}_Y:=\exists_{\pi_Y}\top.
\]
A non-vacuous universal projection is obtained by replacing
\[
  \forall_{\pi_Y}\psi
\]
with
\[
  \operatorname{supp}_Y\wedge \forall_{\pi_Y}\psi.
\]
Its coequational approximation is
\[
  \mathrm{DerInt}
  \left(
    \operatorname{supp}_Y\wedge \forall_{\pi_Y}\psi
  \right).
\]

The right adjoint property in
Theorem~\ref{thm:may-must-coequational-approximation} belongs to the
vacuous universal projection.  For the non-vacuous construction, restrict the
coequational side to the \(Y\)-supported local coequational theories
\[
  \mathbf{Coeq}^{Y\text{-supp}}_{A,B}
  :=
  \left\{
    V\in\mathbf{Coeq}_{A,B}
    \ \middle|\
    iV\leq \operatorname{supp}_Y
  \right\}.
\]
Then for every \(V\in \mathbf{Coeq}^{Y\text{-supp}}_{A,B}\),
\[
  J_Y(V)\leq\psi
  \quad\Longleftrightarrow\quad
  V\leq
  \mathrm{DerInt}
  \left(
    \operatorname{supp}_Y\wedge \forall_{\pi_Y}\psi
  \right).
\]
Thus the non-vacuous construction gives the largest \(Y\)-supported local
coequational theory whose \(Y\)-extension lies in \(\psi\).
\end{remark}

\section{Coequations as Dynamic Invariants}
\label{sec:coequations-as-dynamic-invariants}

Local coequational theories are stable under residual derivatives.  Relative
execution turns residual derivatives into primitive dynamic steps.  Therefore
every local coequational theory induces an invariant predicate on canonical
relative execution: once the residual behaviour lies in the theory, every
primitive step keeps the residual behaviour inside the theory.

Let
\[
  V\hookrightarrow\mathrm{Beh}_0
\]
be a local coequational theory.  Its canonical \(Y\)-extension is
\[
  V_Y:=J_Y(V)=V\times_{B_0}Y.
\]
The primitive canonical box is the predicate transformer
\[
  [m_{\mathrm{Beh},Y}](\varphi)
  :=
  \forall_{s_{\mathrm{Beh},Y}}
  t_{\mathrm{Beh},Y}^*(\varphi).
\]

\begin{theorem}[Coequations induce primitive invariants]
\label{thm:coequations-induce-primitive-invariants}
For every local coequational theory \(V\), one has
\[
  V_Y\leq [m_{\mathrm{Beh},Y}]V_Y.
\]
\end{theorem}

\begin{proof}
\leavevmode
\begin{enumerate}
  \item \textbf{Transpose the desired invariant.}
  By the adjunction
  \[
    s_{\mathrm{Beh},Y}^*\dashv\forall_{s_{\mathrm{Beh},Y}},
  \]
  it is enough to show
  \[
    s_{\mathrm{Beh},Y}^*V_Y
    \leq
    t_{\mathrm{Beh},Y}^*V_Y.
  \]

  \item \textbf{Describe a source element.}
  A generalized element of \(s_{\mathrm{Beh},Y}^*V_Y\) is a triple
  \[
    (a,H,y)
  \]
  with
  \[
    t_A(a)=p_{\mathrm{Beh}}(H),
    \qquad
    q_{\mathrm{Beh}}(H)=r_Y(y),
    \qquad
    H\in V.
  \]

  \item \textbf{Compute its target.}
  Its target is
  \[
    t_{\mathrm{Beh},Y}(a,H,y)
    =
    \left(
      \partial_aH,
      \sigma_Y(\omega_{\mathrm{Beh}}(a,H),y)
    \right).
  \]

  \item \textbf{Use derivative closure.}
  Since \(V\) is derivative-closed,
  \[
    H\in V
    \quad\Longrightarrow\quad
    \partial_aH\in V.
  \]
  Therefore the target lies in
  \[
    V\times_{B_0}Y=V_Y.
  \]

  \item \textbf{Conclude the invariant.}
  Hence
  \[
    s_{\mathrm{Beh},Y}^*V_Y
    \leq
    t_{\mathrm{Beh},Y}^*V_Y.
  \]
  Transposing along
  \(s_{\mathrm{Beh},Y}^*\dashv\forall_{s_{\mathrm{Beh},Y}}\) gives
  \[
    V_Y\leq [m_{\mathrm{Beh},Y}]V_Y.
  \]
\end{enumerate}
\end{proof}

\begin{corollary}[Pullback invariance]
\label{cor:pullback-invariance}
For every Mealy morphism
\[
  \Phi:P:A\rightsquigarrow B,
\]
one has
\[
  \beta_{P,Y}^*V_Y
  \leq
  [m_{\Phi,Y}]\beta_{P,Y}^*V_Y.
\]
If
\[
  P\models V,
\]
that is, if \(\beta_P:P\to\mathrm{Beh}_0\) factors through \(V\), then
\[
  \beta_{P,Y}^*V_Y
  =
  \top_{S_{\Phi,Y}}.
\]
\end{corollary}

\begin{proof}
\leavevmode
\begin{enumerate}
  \item \textbf{Pull back the canonical invariant.}
  Pull back
  \[
    V_Y\leq [m_{\mathrm{Beh},Y}]V_Y
  \]
  along
  \[
    \beta_{P,Y}:S_{\Phi,Y}\to S^{\mathrm{can}}_Y.
  \]

  \item \textbf{Use primitive canonical pullback semantics.}
  By Corollary~\ref{cor:primitive-canonical-pullback-semantics}, the
  pullback of the canonical primitive box is the primitive box for
  \(m_{\Phi,Y}\).  Hence
  \[
    \beta_{P,Y}^*V_Y
    \leq
    [m_{\Phi,Y}]\beta_{P,Y}^*V_Y.
  \]

  \item \textbf{Handle the model case.}
  If \(P\models V\), then \(\beta_P\) factors through \(V\).  Therefore
  \[
    \beta_{P,Y}:S_{\Phi,Y}\to S^{\mathrm{can}}_Y
  \]
  factors through
  \[
    V_Y\hookrightarrow S^{\mathrm{can}}_Y.
  \]
  Hence
  \[
    \beta_{P,Y}^*V_Y
    =
    \top_{S_{\Phi,Y}}.
  \]
\end{enumerate}
\end{proof}

\section{Formula-Generated Coequational Theories and Recognizers}
\label{sec:formula-generated-theories-and-recognizers}

A formula first denotes a predicate on a canonical execution object.  The may
and must coequationalizations extract derivative-stable behaviour
specifications from that predicate.  Applying the transition--output image
construction then gives the algebraic recognizers associated to those
specifications.

Let
\[
  S^{\mathrm{can}}_Y\vdash \varphi\ \mathsf{prop}
\]
be a canonical execution formula with interpretation
\[
  \llbracket\varphi\rrbracket_{\mathrm{can},Y}
  \hookrightarrow
  S^{\mathrm{can}}_Y.
\]

\begin{definition}[Formula-generated coequational theories]
\label{def:formula-generated-coequational-theories}
The \textbf{existential coequational theory generated by \(\varphi\)} is
\[
  \mathrm{Th}^{\exists}_Y(\varphi)
  :=
  \mathrm{May}_Y
  \left(
    \llbracket\varphi\rrbracket_{\mathrm{can},Y}
  \right).
\]
Equivalently,
\[
  \mathrm{Th}^{\exists}_Y(\varphi)
  =
  \mathrm{Der}_0
  \left(
    \exists_{\pi_Y}
    \llbracket\varphi\rrbracket_{\mathrm{can},Y}
  \right).
\]

The \textbf{universal coequational theory generated by \(\varphi\)} is
\[
  \mathrm{Th}^{\forall}_Y(\varphi)
  :=
  \mathrm{Must}_Y
  \left(
    \llbracket\varphi\rrbracket_{\mathrm{can},Y}
  \right).
\]
Equivalently,
\[
  \mathrm{Th}^{\forall}_Y(\varphi)
  =
  \mathrm{DerInt}
  \left(
    \forall_{\pi_Y}
    \llbracket\varphi\rrbracket_{\mathrm{can},Y}
  \right).
\]
\end{definition}

\begin{theorem}[Universal properties of formula-generated theories]
\label{thm:formula-generated-universal-properties}
For every local coequational theory \(V\),
\[
  \mathrm{Th}^{\exists}_Y(\varphi)\leq V
\]
if and only if
\[
  \llbracket\varphi\rrbracket_{\mathrm{can},Y}
  \leq
  J_Y(V).
\]
Furthermore,
\[
  V\leq\mathrm{Th}^{\forall}_Y(\varphi)
\]
if and only if
\[
  J_Y(V)\leq
  \llbracket\varphi\rrbracket_{\mathrm{can},Y}.
\]
\end{theorem}

\begin{proof}
\leavevmode
\begin{enumerate}
  \item \textbf{Substitute the formula interpretation.}
  Apply Theorem~\ref{thm:may-must-coequational-approximation} to
  \[
    \psi=\llbracket\varphi\rrbracket_{\mathrm{can},Y}.
  \]

  \item \textbf{Read off the two adjunctions.}
  The left adjunction gives the existential universal property, and the
  right adjunction gives the universal universal property.
\end{enumerate}
\end{proof}

\begin{proposition}[Monotonicity]
\label{prop:formula-generated-monotonicity}
If
\[
  \llbracket\varphi\rrbracket_{\mathrm{can},Y}
  \leq
  \llbracket\psi\rrbracket_{\mathrm{can},Y},
\]
then
\[
  \mathrm{Th}^{\exists}_Y(\varphi)
  \leq
  \mathrm{Th}^{\exists}_Y(\psi),
\]
and
\[
  \mathrm{Th}^{\forall}_Y(\varphi)
  \leq
  \mathrm{Th}^{\forall}_Y(\psi).
\]
\end{proposition}

\begin{proof}
\leavevmode
\begin{enumerate}
  \item \textbf{Use monotonicity of projection.}
  The functors
  \[
    \exists_{\pi_Y},
    \qquad
    \forall_{\pi_Y}
  \]
  are monotone.

  \item \textbf{Use monotonicity of coequationalization.}
  The functors
  \[
    \mathrm{Der}_0,
    \qquad
    \mathrm{DerInt}
  \]
  are monotone.

  \item \textbf{Compose the monotone maps.}
  Since \(\mathrm{Th}^{\exists}_Y\) and
  \(\mathrm{Th}^{\forall}_Y\) are the corresponding composites, both are
  monotone.
\end{enumerate}
\end{proof}

For a local coequational theory
\[
  V\hookrightarrow \mathrm{Beh}_0,
\]
the restricted canonical behaviour object carries a residual
transition--output action
\[
  \chi_V:A^{\mathrm{op}}\longrightarrow \operatorname{KlEnd}_B(V).
\]
Its identity-on-objects regular image is denoted
\[
  \operatorname{Op}(V).
\]

\begin{definition}[Syntactic transition--output recognizers]
\label{def:syntactic-transition-output-recognizers}
For a canonical execution formula \(\varphi\), define its
\textbf{existential syntactic transition--output recognizer} by
\[
  \operatorname{SynTO}^{\exists}_Y(\varphi)
  :=
  \operatorname{Op}
  \left(
    \mathrm{Th}^{\exists}_Y(\varphi)
  \right),
\]
and define its \textbf{universal syntactic transition--output recognizer} by
\[
  \operatorname{SynTO}^{\forall}_Y(\varphi)
  :=
  \operatorname{Op}
  \left(
    \mathrm{Th}^{\forall}_Y(\varphi)
  \right).
\]
\end{definition}

Thus every canonical execution formula gives a pair of algebraic recognizers
through the pipeline
\[
\begin{tikzcd}[column sep=large]
\varphi
  \arrow[r, mapsto]
&
\llbracket\varphi\rrbracket_{\mathrm{can},Y}
  \arrow[r, mapsto]
&
\mathrm{Th}^{\exists/\forall}_Y(\varphi)
  \arrow[r, mapsto]
&
\operatorname{SynTO}^{\exists/\forall}_Y(\varphi).
\end{tikzcd}
\]

\section{Dynamic Span Logic}
\label{sec:dynamic-span-logic-formal}

A dynamic span
\[
  X\xleftarrow{\ell}U\xrightarrow{r}Y
\]
is read as a possibly nondeterministic program from \(X\) to \(Y\).  A point
of \(U\) is an execution witness.  The map \(\ell\) records the starting
state, and the map \(r\) records the ending state.  The diamond modality asks
for the existence of an execution witness ending in a state satisfying the
postcondition.  The box modality asks that every execution witness from the
current state ends in a state satisfying the postcondition.

\begin{definition}[Dynamic span from Chapter~\ref{ch:relative-mealy-program-categories}]
\label{def:dynamic-span-recalled}
A \textbf{dynamic span} in this appendix means one of the spans to which the
dynamic predicate-transformer semantics of
Chapter~\ref{ch:relative-mealy-program-categories}
has already been applied.  Thus its diamond, and its box when present, are
the predicate transformers
\[
  \langle U\rangle\varphi
  :=
  \exists_{\ell}r^*\varphi,
  \qquad
  [U]\varphi
  :=
  \forall_{\ell}r^*\varphi.
\]
\end{definition}

\begin{proposition}[Adjoint rules for span modalities]
\label{prop:adjoint-rules-span-modalities}
Let
\[
  U=(X\xleftarrow{\ell}U\xrightarrow{r}Y)
\]
be a dynamic span from the earlier span semantics.  For predicates
\[
  \theta\hookrightarrow X,
  \qquad
  \varphi\hookrightarrow Y,
\]
one has
\[
  \langle U\rangle\varphi\leq \theta
  \quad\Longleftrightarrow\quad
  r^*\varphi\leq \ell^*\theta,
\]
and, in the box fragment,
\[
  \theta\leq [U]\varphi
  \quad\Longleftrightarrow\quad
  \ell^*\theta\leq r^*\varphi.
\]
\end{proposition}

\begin{proof}
\leavevmode
\begin{enumerate}
  \item \textbf{Transpose the diamond inequality.}
  Since
  \[
    \langle U\rangle\varphi=\exists_{\ell}r^*\varphi,
  \]
  the adjunction
  \[
    \exists_{\ell}\dashv \ell^*
  \]
  gives
  \[
    \langle U\rangle\varphi\leq \theta
    \quad\Longleftrightarrow\quad
    r^*\varphi\leq \ell^*\theta.
  \]

  \item \textbf{Transpose the box inequality.}
  In the box fragment,
  \[
    [U]\varphi=\forall_{\ell}r^*\varphi.
  \]
  The adjunction
  \[
    \ell^*\dashv\forall_{\ell}
  \]
  gives
  \[
    \theta\leq [U]\varphi
    \quad\Longleftrightarrow\quad
    \ell^*\theta\leq r^*\varphi.
  \]
\end{enumerate}
\end{proof}

\begin{definition}[Tests]
\label{def:dynamic-tests}
For a predicate
\[
  \theta\hookrightarrow X
\]
in a dynamic context where the test program was constructed in
Chapter~\ref{ch:relative-mealy-program-categories},
the \textbf{test program}
\[
  \theta?:X\rightsquigarrow X
\]
is the span
\[
  X\xleftarrow{\theta}\Theta\xrightarrow{\theta}X,
\]
where \(\Theta\hookrightarrow X\) represents \(\theta\).
\end{definition}

\begin{proposition}[Test laws]
\label{prop:test-laws}
For every predicate
\[
  \varphi\hookrightarrow X,
\]
one has
\[
  \langle\theta?\rangle\varphi
  =
  \theta\wedge\varphi.
\]
In the Heyting box fragment,
\[
  [\theta?]\varphi
  =
  \theta\Rightarrow\varphi.
\]
\end{proposition}

\begin{proof}
\leavevmode
\begin{enumerate}
  \item \textbf{Compute the diamond.}
  Since both legs of the test span are the inclusion
  \[
    \theta:\Theta\hookrightarrow X,
  \]
  the diamond is
  \[
    \exists_{\theta}\theta^*\varphi.
  \]
  This is the image of the pullback of \(\varphi\) to \(\Theta\), hence the
  meet
  \[
    \theta\wedge\varphi.
  \]

  \item \textbf{Compute the box.}
  The box is
  \[
    \forall_{\theta}\theta^*\varphi.
  \]
  This is characterized by the right adjoint to pullback along
  \(\theta\).  In the Heyting subobject semantics, this right adjoint is
  \[
    \theta\Rightarrow\varphi.
  \]
\end{enumerate}
\end{proof}

\begin{definition}[Finite choice]
\label{def:finite-choice-subprograms}
For subprograms
\[
  U,V\hookrightarrow M
\]
of a common dynamic span
\[
  X\xleftarrow{\ell}M\xrightarrow{r}Y
\]
in the finite-choice fragment of
Chapter~\ref{ch:relative-mealy-program-categories},
their \textbf{finite choice} is the join
\[
  U\vee V\hookrightarrow M.
\]
\end{definition}

\begin{proposition}[Finite-choice laws]
\label{prop:finite-choice-laws}
For finite-choice subprograms of a common dynamic span, one has
\[
  \langle U\vee V\rangle\varphi
  =
  \langle U\rangle\varphi
  \vee
  \langle V\rangle\varphi.
\]
In the box fragments where the corresponding universal decomposition was
constructed, one has
\[
  [U\vee V]\varphi
  =
  [U]\varphi
  \wedge
  [V]\varphi.
\]
In particular, the box law holds for homwise coproduct choice in the modal
feedback bases of Chapter~\ref{ch:feedback-trace}.
\end{proposition}

\begin{proof}
\leavevmode
\begin{enumerate}
  \item \textbf{Use finite-choice semantics for diamonds.}
  The diamond law is the preservation of finite joins by existential image
  along the common source map:
  \[
    \langle U\vee V\rangle\varphi
    =
    \langle U\rangle\varphi
    \vee
    \langle V\rangle\varphi.
  \]

  \item \textbf{Use finite-choice semantics for boxes.}
  In the fragments where boxes for finite choice were constructed, the
  corresponding universal decomposition gives
  \[
    [U\vee V]\varphi
    =
    [U]\varphi
    \wedge
    [V]\varphi.
  \]
\end{enumerate}
\end{proof}

\begin{definition}[Dynamic formula-generated theories]
\label{def:dynamic-formula-generated-theories}
For a dynamic formula
\[
  S^{\mathrm{can}}_Y\vdash\varphi\ \mathsf{prop},
\]
define its \textbf{existential dynamic formula-generated theory} by
\[
  \mathrm{Th}^{\exists,\mathrm{dyn}}_Y(\varphi)
  :=
  \mathrm{Th}^{\exists}_Y(\varphi),
\]
and define its \textbf{universal dynamic formula-generated theory} by
\[
  \mathrm{Th}^{\forall,\mathrm{dyn}}_Y(\varphi)
  :=
  \mathrm{Th}^{\forall}_Y(\varphi).
\]
\end{definition}

Thus dynamic span logic contributes canonical execution predicates, and the
coequational semantics of those predicates is the may--must construction of
Theorem~\ref{thm:may-must-coequational-approximation}.

\section{Guarded Response-Profile Logic}
\label{sec:guarded-response-profile-logic-formal}

The response-profile language inspects the labelled data of one execution
step.  A guarded profile can mention the input arrow, the output-comparison
arrow, the downstream transition, and the next state.  Flattening identifies
this polynomial description with an ordinary dynamic span whose apex is the
object of response positions satisfying the same guards.

This section uses the relative response polynomial
\[
  M_{A,B,Y}:\mathcal E/S\to \mathcal E/S
\]
constructed in Chapter~\ref{ch:polynomial-relative-mealy-actions}.  For each
state object \(S\) over the relative interface, write
\[
  M_{A,B,Y}(S)
\]
for its response object.

\begin{definition}[Response positions]
\label{def:response-positions}
For a state object \(S\), let
\[
  \operatorname{Pos}_{A,B,Y}(S)
\]
be the \textbf{position object} of the response polynomial over \(S\).  It is
equipped with a position projection
\[
  \partial_S:
  \operatorname{Pos}_{A,B,Y}(S)
  \longrightarrow
  M_{A,B,Y}(S),
\]
and structural maps
\[
  \mathrm{in}:
  \operatorname{Pos}_{A,B,Y}(S)\to A_1,
\]
\[
  \mathrm{out}:
  \operatorname{Pos}_{A,B,Y}(S)\to B_1,
\]
\[
  \mathrm{down}:
  \operatorname{Pos}_{A,B,Y}(S)\to
  \left(\int_B Y\right)_1,
\]
and
\[
  \mathrm{next}:
  \operatorname{Pos}_{A,B,Y}(S)\to S.
\]
\end{definition}

Let
\[
  R\hookrightarrow A_1,
  \qquad
  O\hookrightarrow B_1,
  \qquad
  T\hookrightarrow \left(\int_B Y\right)_1
\]
be predicates.  Let
\[
  G^S_{R/O/T}\hookrightarrow \operatorname{Pos}_{A,B,Y}(S)
\]
be the response-position predicate
\[
  G^S_{R/O/T}
  :=
  \mathrm{in}^*R
  \wedge
  \mathrm{out}^*O
  \wedge
  \mathrm{down}^*T.
\]
Write
\[
  \partial^G_S:G^S_{R/O/T}\to M_{A,B,Y}(S),
  \qquad
  \mathrm{next}^G:G^S_{R/O/T}\to S
\]
for the restrictions of \(\partial_S\) and \(\mathrm{next}\).

For a predicate
\[
  \varphi\hookrightarrow S,
\]
the guarded existential profile lifting constructed in
Chapter~\ref{ch:polynomial-relative-mealy-actions} is
\[
  \lambda_{R/O/T}^S(\varphi)
  \hookrightarrow
  M_{A,B,Y}(S),
\]
given by
\[
  \lambda_{R/O/T}^S(\varphi)
  :=
  \exists_{\partial^G_S}
  \left(
    (\mathrm{next}^G)^*\varphi
  \right).
\]
The weak guarded universal profile lifting is the corresponding box lifting
that agrees with the flattened guarded span box; write it as
\[
  \rho^{\mathrm{wk},S}_{R/O/T}(\varphi).
\]

For a Mealy morphism
\[
  \Phi:P:A\rightsquigarrow B
\]
and a right \(B\)-action \(Y\), let
\[
  c_{\Phi,Y}:S_{\Phi,Y}\to M_{A,B,Y}(S_{\Phi,Y})
\]
be the induced response coalgebra.  The canonical response coalgebra is
\[
  c_{\mathrm{Beh},Y}:
  S^{\mathrm{can}}_Y
  \to
  M_{A,B,Y}(S^{\mathrm{can}}_Y).
\]

\begin{definition}[Guarded profile modalities]
\label{def:guarded-profile-modalities}
Let
\[
  \theta\hookrightarrow S_{\Phi,Y}
\]
be a state guard.  The \textbf{guarded profile diamond} is the modality
constructed in Chapter~\ref{ch:polynomial-relative-mealy-actions}; its
interpretation is
\[
  \llbracket
    \langle\theta;R/O/T\rangle\varphi
  \rrbracket_{\Phi,Y}
  :=
  \theta
  \wedge
  c_{\Phi,Y}^*
  \lambda_{R/O/T}^{S_{\Phi,Y}}
  \left(
    \llbracket\varphi\rrbracket_{\Phi,Y}
  \right).
\]

The \textbf{weak guarded profile box} is the box that agrees with the box of
the flattened guarded span:
\[
  \llbracket
    [\theta;R/O/T]_{\mathrm{wk}}\varphi
  \rrbracket_{\Phi,Y}
  :=
  \theta
  \Rightarrow
  c_{\Phi,Y}^*
  \rho^{\mathrm{wk},S_{\Phi,Y}}_{R/O/T}
  \left(
    \llbracket\varphi\rrbracket_{\Phi,Y}
  \right).
\]

The \textbf{strong guarded profile box} is the stronger profile-level
predicate transformer of Chapter~\ref{ch:polynomial-relative-mealy-actions}.
It is first interpreted as a predicate on the canonical relative execution
object \(S^{\mathrm{can}}_Y\), via the canonical response coalgebra
\(c_{\mathrm{Beh},Y}\).  Its coequational semantics is then obtained by
applying \(\mathrm{May}_Y\) or \(\mathrm{Must}_Y\) to that canonical
predicate.
\end{definition}

\begin{definition}[Flattening of a guarded profile]
\label{def:flattening-guarded-profile}
The \textbf{flattened span program} associated to
\[
  \langle\theta;R/O/T\rangle
\]
is the guarded subspan of the primitive execution span whose apex is the
subobject of response positions satisfying
\[
  \theta,
  \qquad
  R,
  \qquad
  O,
  \qquad
  T.
\]
It is denoted
\[
  \langle\theta;R/O/T\rangle_{\mathrm{span}}.
\]
The flattened span program associated to
\[
  [\theta;R/O/T]_{\mathrm{wk}}
\]
uses the same guarded response-position apex and the box predicate
transformer of that span.
\end{definition}

The comparison between profile semantics and span semantics is represented
by
\[
\begin{tikzcd}[column sep=large,row sep=large]
G^S_{R/O/T}
  \arrow[r, hook]
  \arrow[d, "\partial^G_S"']
  \arrow[dr, phantom, very near start, "\lrcorner"]
&
\operatorname{Pos}_{A,B,Y}(S)
  \arrow[d, "\partial_S"]
  \arrow[r, "\mathrm{next}"]
&
S
\\
M_{A,B,Y}(S)
  &
M_{A,B,Y}(S)
  \arrow[l, equal]
&
\end{tikzcd}
\]
together with the coalgebra map
\[
  c_{\Phi,Y}:S_{\Phi,Y}\to M_{A,B,Y}(S_{\Phi,Y}).
\]
Pulling guarded profile predicates back along \(c_{\Phi,Y}\) produces the
same predicate transformer as the guarded subspan obtained by flattening.

\begin{theorem}[Profile flattening]
\label{thm:profile-flattening}
For every guarded profile diamond formula
\[
  \rho=\langle\theta;R/O/T\rangle\varphi,
\]
one has
\[
  \llbracket \rho\rrbracket_{\Phi,Y}
  =
  \llbracket \rho^\flat\rrbracket_{\Phi,Y},
\]
where \(\rho^\flat\) is the corresponding flattened dynamic-span formula.
The same equality holds for weak guarded profile boxes.  Strong profile boxes
are interpreted as canonical predicates and then coequationalized by the same
may--must construction.  In particular,
\[
  \llbracket \rho\rrbracket_{\mathrm{can},Y}
  =
  \llbracket \rho^\flat\rrbracket_{\mathrm{can},Y}
\]
for diamonds and weak boxes.
\end{theorem}

\begin{proof}
\leavevmode
\begin{enumerate}
  \item \textbf{Identify the guarded response-position object.}
  The response-profile diamond interpretation is the pullback along
  \[
    c_{\Phi,Y}
  \]
  of an existential projection along the guarded position projection
  \[
    \partial^G_S:G^S_{R/O/T}\to M_{A,B,Y}(S).
  \]

  \item \textbf{Compare with the flattened span.}
  Flattening identifies \(G^S_{R/O/T}\) with the apex of a guarded subspan of
  the primitive execution span.  The source and target maps of that subspan
  are the current-state and continuation maps induced by the response
  position.

  \item \textbf{Conclude equality for diamonds.}
  Under this identification, both semantics compute the same existential
  predicate transformer.  Therefore the profile diamond and its flattened
  dynamic-span diamond have equal interpretations.

  \item \textbf{Conclude equality for weak boxes.}
  The weak box interpretation is the corresponding universal projection for
  the same guarded response-position object.  Hence it agrees with the box of
  the flattened guarded span.

  \item \textbf{Handle the canonical case.}
  The same argument applies with
  \[
    S=S^{\mathrm{can}}_Y.
  \]
  Thus canonical interpretations agree for diamonds and weak boxes.

  \item \textbf{Record the strong-box case.}
  Strong profile boxes are not identified with flattened span boxes here.
  They are treated as canonical predicates and then passed through the same
  may--must coequationalization.
\end{enumerate}
\end{proof}

\begin{corollary}[Profile coequationalization agrees with flattening]
\label{cor:profile-coequationalization-flattening}
For every guarded profile diamond formula \(\rho\),
\[
  \mathrm{Th}^{\exists}_Y(\rho)
  =
  \mathrm{Th}^{\exists}_Y(\rho^\flat).
\]
For every weak guarded profile box formula \(\rho\),
\[
  \mathrm{Th}^{\forall}_Y(\rho)
  =
  \mathrm{Th}^{\forall}_Y(\rho^\flat).
\]
\end{corollary}

\begin{proof}
\leavevmode
\begin{enumerate}
  \item \textbf{Use canonical equality.}
  By Theorem~\ref{thm:profile-flattening}, the canonical interpretations
  agree.

  \item \textbf{Apply coequationalization.}
  Applying
  \[
    \mathrm{May}_Y
    \qquad
    \text{and}
    \qquad
    \mathrm{Must}_Y
  \]
  gives the displayed equalities of generated coequational theories.
\end{enumerate}
\end{proof}

\section{Cubical Execution Logic and Sequential Projections}
\label{sec:cubical-execution-logic-formal}

Cubical execution predicates may observe higher-dimensional execution cells.
Such predicates live on cubical execution objects rather than on ordinary
residual behaviour objects.  The construction in this section does not define
cubical coequations.  It extracts ordinary residual-behaviour coequations
from cubical observations by applying one of the sequential behaviour
projections already present in the cubical semantics.

Let \(F\) be a finite tensor context of evaluated Mealy executions
\[
  (\Phi_i,Y_i)_{i\in I}.
\]
Write
\[
  S_F:=\prod_{i\in I}S_{\Phi_i,Y_i}
\]
for the product state object.  For a Boolean interval
\[
  \alpha\leq\beta
\]
in \(2^I\), write
\[
  E^F_{\alpha,\beta}
\]
for the corresponding interval-cubical execution object.  Its lower and upper
endpoint maps are
\[
  d^\alpha_{\alpha,\beta},
  \qquad
  d^\beta_{\alpha,\beta}.
\]
The canonical interval-cubical execution object is denoted
\[
  E^{\mathrm{can},F}_{\alpha,\beta}.
\]

Residual semantics induces comparison maps
\[
  \beta^\square_{F,\alpha,\beta}:
  E^F_{\alpha,\beta}
  \longrightarrow
  E^{\mathrm{can},F}_{\alpha,\beta}.
\]
For a cubical context \(\Gamma\), write
\[
  \beta^\Gamma_F:
  \llbracket\Gamma\rrbracket_F
  \longrightarrow
  \llbracket\Gamma\rrbracket_{\mathrm{can},F}
\]
for the induced comparison map.

\begin{definition}[Cell diamond]
\label{def:cell-diamond}
Let
\[
  U\hookrightarrow E^F_{\alpha,\beta}
\]
be a guarded interval subobject from the dynamic-CETT fragment.  The
\textbf{cell diamond} is interpreted by
\[
  \llbracket
    \langle U\rangle_{\alpha,\beta}\varphi
  \rrbracket_F
  :=
  \exists_{d^\alpha_{\alpha,\beta}|_U}
  \left(
    (d^\beta_{\alpha,\beta}|_U)^*
    \llbracket\varphi\rrbracket_F
  \right).
\]
\end{definition}

\begin{theorem}[Cubical truth is pulled back from canonical cubical execution]
\label{thm:cubical-truth-pullback}
Let
\[
  \Gamma\vdash\chi\ \mathsf{prop}
\]
be a CETT or dynamic-CETT formula from the fragments constructed in
Chapter~\ref{ch:cubical-mealy-execution-concurrency-types}.  If its atomic
predicates and guarded cell subobjects are canonical, then
\[
  \llbracket\chi\rrbracket_F
  =
  (\beta^\Gamma_F)^*
  \llbracket\chi\rrbracket_{\mathrm{can},F}.
\]
For formulas involving finite iteration, this statement is read in the
path-star fragment already constructed using the finite-path semantics of
Chapter~\ref{ch:feedback-trace}.
\end{theorem}

\begin{proof}
\leavevmode
\begin{enumerate}
  \item \textbf{Handle the finite-limit fragment.}
  The finite-limit part of CETT is interpreted by products, pullbacks, face
  maps, endpoint maps, and profile maps.  These constructions are natural
  under residual semantics.  Hence the claim holds for the finite-limit
  fragment.

  \item \textbf{Set up the guarded cell square.}
  For a cell diamond, canonicality of the guard gives a pullback square
  \[
  \begin{tikzcd}[column sep=large,row sep=large]
  U_F
    \arrow[r]
    \arrow[d, "d^\alpha|_{U_F}"']
  &
  U_{\mathrm{can}}
    \arrow[d, "d^\alpha|_{U_{\mathrm{can}}}"]
  \\
  \llbracket\Gamma\rrbracket_F
    \arrow[r, "\beta^\Gamma_F"']
  &
  \llbracket\Gamma\rrbracket_{\mathrm{can},F}
    \arrow[ul, phantom, very near start, "\lrcorner"]
  \end{tikzcd}
  \]
  and the upper endpoint maps commute with the same comparison:
  \[
  \begin{tikzcd}[column sep=large]
  U_F
    \arrow[r]
    \arrow[d, "d^\beta|_{U_F}"']
  &
  U_{\mathrm{can}}
    \arrow[d, "d^\beta|_{U_{\mathrm{can}}}"]
  \\
  \llbracket\Gamma\rrbracket_F
    \arrow[r, "\beta^\Gamma_F"']
  &
  \llbracket\Gamma\rrbracket_{\mathrm{can},F}.
  \end{tikzcd}
  \]

  \item \textbf{Apply Beck--Chevalley.}
  Beck--Chevalley for the lower endpoint map gives compatibility of the
  existential predicate transformer with residual pullback.  The box case,
  when included in the formula syntax, uses the right-adjoint
  Beck--Chevalley law.

  \item \textbf{Handle finite iteration.}
  For finite iteration modalities, the result is the corresponding
  path-star/fixed-point compatibility from the finite-path semantics already
  used by the dynamic cubical fragment.

  \item \textbf{Conclude by induction.}
  The result follows by induction on formulas.
\end{enumerate}
\end{proof}

Cubical formulas produce predicates on cubical execution objects.  To obtain
ordinary local coequational theories, one uses a sequential behaviour
projection already supplied by the cubical observation.

\begin{definition}[Sequential behaviour projection supplied by a cubical observation]
\label{def:sequential-behaviour-projection}
A \textbf{sequential behaviour projection} is one of the behaviour
projections arising from the cubical execution constructions of
Chapter~\ref{ch:cubical-mealy-execution-concurrency-types}.  For example,
in a single-component interval, the lower endpoint residual-behaviour
projection gives
\[
  q_{\alpha,\beta}:
  E^{\mathrm{can},F}_{\alpha,\beta}
  \longrightarrow
  \mathrm{Beh}_{A^\sharp,B^\sharp,0}.
\]
This appendix introduces no additional cubical behaviour object and no
additional cubical-to-sequential projection.
\end{definition}

\begin{definition}[Cubical sequential coequational projections]
\label{def:cubical-sequential-coequational-projections}
Let
\[
  \chi\hookrightarrow E^{\mathrm{can},F}_{\alpha,\beta}
\]
be a canonical cubical predicate, and let
\[
  q_{\alpha,\beta}:
  E^{\mathrm{can},F}_{\alpha,\beta}
  \to
  \mathrm{Beh}_{A^\sharp,B^\sharp,0}
\]
be a sequential behaviour projection supplied by the cubical semantics.
Define the \textbf{existential cubical sequential coequational projection} by
\[
  \mathrm{Th}^{\exists,\square}_{\alpha,\beta}(\chi)
  :=
  \mathrm{Der}^{A^\sharp,B^\sharp}_0
  \left(
    \exists_{q_{\alpha,\beta}}\chi
  \right),
\]
and define the \textbf{universal cubical sequential coequational projection}
by
\[
  \mathrm{Th}^{\forall,\square}_{\alpha,\beta}(\chi)
  :=
  \mathrm{DerInt}^{A^\sharp,B^\sharp}
  \left(
    \forall_{q_{\alpha,\beta}}\chi
  \right).
\]
These are local coequational theories in
\[
  \mathbf{Coeq}_{A^\sharp,B^\sharp}.
\]
\end{definition}

\begin{remark}[Ordinary rather than cubical coequations]
\label{rem:ordinary-rather-than-cubical-coequations}
The construction above produces ordinary residual-behaviour projections of
cubical observations.  A full non-interleaving coequational theory would
require a cubical behaviour object, for example an object
\[
  \mathrm{Beh}^{\square}_{A,B},
\]
together with an appropriate notion of cubical derivative-closed subobject.
The present appendix uses the ordinary coequational semantics already
constructed for residual behaviours.
\end{remark}

\section{Finite Paths and Feedback Trace}
\label{sec:finite-paths-feedback-trace-formal}

Finite-path semantics internalizes repeated execution.  The diamond
\[
  \langle p^\star\rangle\varphi
\]
says that some finite number of \(p\)-steps reaches \(\varphi\).  The box
\[
  [p^\star]\varphi
\]
says that all finite \(p\)-executions remain within the specified
postcondition, in the sense encoded by the finite-path semantics of
Chapter~\ref{ch:feedback-trace}.  Feedback trace
adds a hidden interface: an execution may move from the visible part into the
hidden part, iterate there for finitely many hidden steps, and then return to
the visible part.

In this section we use the modal feedback base of
Chapter~\ref{ch:feedback-trace}.  Thus finite path
spans, feedback trace spans, fixed-point interpretations of the corresponding
predicate transformers, and the associated Beck--Chevalley laws are exactly
the structures already used there.

Let
\[
  p:X\rightsquigarrow X
\]
be an endospan in the finite-path fragment.  Its finite-path span is denoted
\[
  p^\star:X\rightsquigarrow X.
\]
The diamond of finite iteration is interpreted by the least fixed point
\[
  \langle p^\star\rangle\varphi
  =
  \mu Z.\,(\varphi\vee \langle p\rangle Z).
\]
The box of finite iteration is interpreted by the greatest fixed point
\[
  [p^\star]\varphi
  =
  \nu Z.\,(\varphi\wedge [p]Z).
\]

Let
\[
  R:X+H\rightsquigarrow Y+H
\]
be a feedback span.  Its block decomposition is written
\[
  R=
  \begin{pmatrix}
    R_{XY} & R_{XH}\\
    R_{HY} & R_{HH}
  \end{pmatrix}.
\]
The feedback trace is
\[
  \operatorname{Tr}^H(R)
  =
  R_{XY}
  \boxplus
  R_{HY}\circ \operatorname{Path}(R_{HH})\circ R_{XH}.
\]

When the box fragment of the feedback semantics is available, the box
\[
  [\operatorname{Tr}^H(R)]
\]
denotes the universal predicate transformer of the traced feedback span
constructed in
Chapter~\ref{ch:feedback-trace}.

\begin{proposition}[Trace diamond]
\label{prop:trace-diamond}
For a postcondition
\[
  \varphi\hookrightarrow Y,
\]
one has
\[
  \langle\operatorname{Tr}^H(R)\rangle\varphi
  =
  \langle R_{XY}\rangle\varphi
  \vee
  \langle R_{XH}\rangle
  \mu Z.
  \left(
    \langle R_{HY}\rangle\varphi
    \vee
    \langle R_{HH}\rangle Z
  \right).
\]
\end{proposition}

\begin{proof}
\leavevmode
\begin{enumerate}
  \item \textbf{Decompose the trace span.}
  The trace span decomposes as the finite choice of the direct visible branch
  \[
    R_{XY}
  \]
  and the branch that enters the hidden interface through \(R_{XH}\), iterates
  inside \(H\) along \(R_{HH}\), and exits through \(R_{HY}\).

  \item \textbf{Apply finite choice.}
  The diamond of a finite choice is the join of the diamonds of the chosen
  branches.

  \item \textbf{Apply finite-path semantics.}
  The hidden iteration is interpreted by the least fixed point
  \[
    \mu Z.
    \left(
      \langle R_{HY}\rangle\varphi
      \vee
      \langle R_{HH}\rangle Z
    \right).
  \]

  \item \textbf{Assemble the formula.}
  Combining the direct branch and hidden branch gives the displayed
  expression.
\end{enumerate}
\end{proof}

Let
\[
  S=S^{\mathrm{can}}_Y.
\]
Let
\[
  U:S\rightsquigarrow S
\]
be a canonical relative execution sub-endospan.  Define
\[
  L_U:S+S\rightsquigarrow S+S
\]
by
\[
  L_U
  =
  \begin{pmatrix}
    0_{S,S} & 1_S\\
    1_S & U
  \end{pmatrix},
\]
where the first copy of \(S\) is visible and the second copy is hidden.

\begin{theorem}[Path-star as feedback trace]
\label{thm:path-star-as-feedback-trace}
There is an isomorphism of spans
\[
  \operatorname{Tr}^S(L_U)
  \cong
  \operatorname{Path}(U).
\]
Consequently, for every postcondition
\[
  \varphi\hookrightarrow S,
\]
one has
\[
  \mathrm{Th}^{\exists}_Y
  \left(
    \langle U^\star\rangle\varphi
  \right)
  =
  \mathrm{Th}^{\exists}_Y
  \left(
    \langle\operatorname{Tr}^S(L_U)\rangle\varphi
  \right),
\]
and
\[
  \mathrm{Th}^{\forall}_Y
  \left(
    [U^\star]\varphi
  \right)
  =
  \mathrm{Th}^{\forall}_Y
  \left(
    [\operatorname{Tr}^S(L_U)]\varphi
  \right).
\]
\end{theorem}

\begin{proof}
\leavevmode
\begin{enumerate}
  \item \textbf{Compute the trace blockwise.}
  For the matrix \(L_U\), the visible-to-visible block is empty, the
  visible-to-hidden block is the identity, the hidden-to-visible block is the
  identity, and the hidden-to-hidden block is \(U\).  Hence the trace formula
  gives
  \[
    \operatorname{Tr}^S(L_U)
    =
    0_{S,S}
    \boxplus
    1_S\circ\operatorname{Path}(U)\circ 1_S.
  \]

  \item \textbf{Cancel the identities and empty branch.}
  The right-hand side is canonically isomorphic to
  \[
    \operatorname{Path}(U).
  \]

  \item \textbf{Pass to predicate transformers.}
  Isomorphic spans induce equal predicate transformers.  Therefore
  \[
    \langle U^\star\rangle\varphi
    =
    \langle\operatorname{Tr}^S(L_U)\rangle\varphi
  \]
  and
  \[
    [U^\star]\varphi
    =
    [\operatorname{Tr}^S(L_U)]\varphi
  \]
  in the corresponding box fragment.

  \item \textbf{Apply coequationalization.}
  Applying
  \[
    \mathrm{May}_Y
    \qquad
    \text{and}
    \qquad
    \mathrm{Must}_Y
  \]
  gives the displayed equalities of coequational theories.
\end{enumerate}
\end{proof}

\section{Stability under Downstream Change and Simulation}
\label{sec:stability-downstream-simulation-formal}

The constructions are functorial in two directions.  Changing the downstream
action changes the environment in which output comparisons are evaluated.  A
morphism of right actions therefore gives substitution of predicates along
the induced map of execution objects.  A Mealy simulation changes the program
being observed.  Since relative program categories are contravariant in
simulations, simulation induces predicate substitution in the opposite
direction.

\subsection{Downstream Change}
\label{subsec:downstream-change-formal}

Let
\[
  g:Y\to Y'
\]
be a morphism of right \(B\)-actions.  It induces maps
\[
  1_P\times_{B_0} g:
  S_{\Phi,Y}
  \longrightarrow
  S_{\Phi,Y'},
\]
and
\[
  h:=1_{\mathrm{Beh}}\times_{B_0} g:
  S^{\mathrm{can}}_Y
  \longrightarrow
  S^{\mathrm{can}}_{Y'}.
\]
The canonical projections satisfy
\[
  \pi_Y=\pi_{Y'}h.
\]

The downstream comparison triangle is
\[
\begin{tikzcd}[column sep=large,row sep=large]
S^{\mathrm{can}}_Y
  \arrow[rr, "h"]
  \arrow[dr, "\pi_Y"']
&&
S^{\mathrm{can}}_{Y'}
  \arrow[dl, "\pi_{Y'}"]
\\
&
\mathrm{Beh}_0 .
\end{tikzcd}
\]

\begin{proposition}[Downstream substitution]
\label{prop:downstream-substitution}
For every formula
\[
  S_{\Phi,Y'}\vdash\varphi\ \mathsf{prop}
\]
in the dynamic, guarded, cubical, or trace fragment under consideration,
there is a substituted formula
\[
  S_{\Phi,Y}\vdash (1_P\times_{B_0} g)^*\varphi\ \mathsf{prop}.
\]
Moreover,
\[
  \varphi\leq\psi
  \quad\Longrightarrow\quad
  (1_P\times_{B_0} g)^*\varphi
  \leq
  (1_P\times_{B_0} g)^*\psi.
\]
\end{proposition}

\begin{proof}
\leavevmode
\begin{enumerate}
  \item \textbf{Handle atomic predicates.}
  Atomic predicates are pulled back along
  \[
    1_P\times_{B_0} g.
  \]

  \item \textbf{Handle propositional structure.}
  Pullback preserves the propositional structure of the fragment.

  \item \textbf{Handle modalities.}
  For modalities, equivariance of \(g\) gives the corresponding morphisms of
  relative execution spans.  The relevant source squares are the
  Beck--Chevalley squares already used in the earlier downstream-stability
  results.

  \item \textbf{Handle guarded, cubical, path, and trace constructors.}
  These constructors are stable by the same functoriality of their execution
  objects.

  \item \textbf{Conclude by induction.}
  The substituted formula exists by induction on formula formation, and
  monotonicity follows because pullback is monotone.
\end{enumerate}
\end{proof}

\begin{proposition}[Downstream comparison of coequational projections]
\label{prop:downstream-comparison-coequational-projections}
Let
\[
  \psi'\hookrightarrow S^{\mathrm{can}}_{Y'}
\]
be a canonical execution predicate.  Then
\[
  \mathrm{May}_Y(h^*\psi')
  \leq
  \mathrm{May}_{Y'}(\psi').
\]
Furthermore,
\[
  \mathrm{Must}_{Y'}(\psi')
  \leq
  \mathrm{Must}_Y(h^*\psi').
\]
If the triangle
\[
  \pi_Y=\pi_{Y'}h
\]
is exact for both existential and universal projection in the sense already
used by the downstream-stability theorem of the relevant fragment, then the
corresponding inequalities are equalities.
\end{proposition}

\begin{proof}
\leavevmode
\begin{enumerate}
  \item \textbf{Compare existential projections.}
  Since
  \[
    \pi_Y=\pi_{Y'}h,
  \]
  the unit of
  \[
    \exists_{\pi_{Y'}}\dashv \pi_{Y'}^*
  \]
  gives
  \[
    \psi'\leq
    \pi_{Y'}^*\exists_{\pi_{Y'}}\psi'.
  \]
  Pulling back along \(h\) gives
  \[
    h^*\psi'
    \leq
    \pi_Y^*\exists_{\pi_{Y'}}\psi'.
  \]
  By the adjunction
  \[
    \exists_{\pi_Y}\dashv \pi_Y^*,
  \]
  this yields
  \[
    \exists_{\pi_Y}h^*\psi'
    \leq
    \exists_{\pi_{Y'}}\psi'.
  \]
  Applying the monotone functor \(\mathrm{Der}_0\) gives
  \[
    \mathrm{May}_Y(h^*\psi')
    \leq
    \mathrm{May}_{Y'}(\psi').
  \]

  \item \textbf{Compare universal projections.}
  The counit of
  \[
    \pi_{Y'}^*\dashv\forall_{\pi_{Y'}}
  \]
  gives
  \[
    \pi_{Y'}^*\forall_{\pi_{Y'}}\psi'
    \leq
    \psi'.
  \]
  Pulling back along \(h\) gives
  \[
    \pi_Y^*\forall_{\pi_{Y'}}\psi'
    \leq
    h^*\psi'.
  \]
  By the adjunction
  \[
    \pi_Y^*\dashv\forall_{\pi_Y},
  \]
  this yields
  \[
    \forall_{\pi_{Y'}}\psi'
    \leq
    \forall_{\pi_Y}h^*\psi'.
  \]
  Applying the monotone functor \(\mathrm{DerInt}\) gives
  \[
    \mathrm{Must}_{Y'}(\psi')
    \leq
    \mathrm{Must}_Y(h^*\psi').
  \]

  \item \textbf{Identify the equality case.}
  If the projection triangle is exact for both existential and universal
  projection in the imported downstream-stability theorem, then the two
  comparison inequalities before applying \(\mathrm{Der}_0\) and
  \(\mathrm{DerInt}\) are equalities.  Hence the corresponding
  coequational comparisons are equalities.
\end{enumerate}
\end{proof}

\subsection{Simulation Substitution}
\label{subsec:simulation-substitution-formal}

Let
\[
  \Theta:\Phi\Rightarrow\Psi
\]
be a Mealy simulation.  Relative program categories are contravariant in
Mealy simulations.  Hence \(\Theta\) induces a functor
\[
  H_{\Theta,Y}:
  \operatorname{Prog}_{\Psi,Y}
  \longrightarrow
  \operatorname{Prog}_{\Phi,Y}.
\]
On state objects, write the associated comparison as
\[
  h_{\Theta,Y}:
  S_{\Psi,Y}\longrightarrow S_{\Phi,Y}.
\]
It induces predicate substitution
\[
  h_{\Theta,Y}^*:
  \operatorname{Pred}(S_{\Phi,Y})
  \longrightarrow
  \operatorname{Pred}(S_{\Psi,Y}).
\]

\begin{proposition}[Simulation substitution]
\label{prop:simulation-substitution-logical}
For every formula
\[
  S_{\Phi,Y}\vdash\varphi\ \mathsf{prop}
\]
in the fragment under consideration, there is a substituted formula
\[
  S_{\Psi,Y}\vdash h_{\Theta,Y}^*\varphi\ \mathsf{prop}.
\]
Moreover,
\[
  \varphi\leq\psi
  \quad\Longrightarrow\quad
  h_{\Theta,Y}^*\varphi
  \leq
  h_{\Theta,Y}^*\psi.
\]
\end{proposition}

\begin{proof}
\leavevmode
\begin{enumerate}
  \item \textbf{Handle atomic predicates.}
  Atomic predicates are pulled back along
  \[
    h_{\Theta,Y}.
  \]

  \item \textbf{Handle propositional structure.}
  Pullback preserves the propositional connectives available in the fragment.

  \item \textbf{Transport dynamic modalities.}
  For dynamic modalities, subprograms are transported along
  \[
    H_{\Theta,Y}.
  \]
  The corresponding span comparison squares are the Beck--Chevalley squares
  from the simulation-stability theorem for relative program categories.

  \item \textbf{Transport higher constructors.}
  The guarded response-profile, cubical, path, and trace cases use the
  functoriality of the corresponding execution objects under Mealy
  simulations.

  \item \textbf{Conclude by induction.}
  The substituted formula exists by induction on formula formation, and
  monotonicity follows from monotonicity of pullback.
\end{enumerate}
\end{proof}

\begin{remark}[Recognizer variance]
\label{rem:recognizer-variance}
The comparison of syntactic transition--output recognizers is obtained by
applying the operation-image construction
\[
  \operatorname{Op}
\]
to the comparison of generated coequational theories.  Its variance is the
variance established for transition--output images in
Chapter~\ref{ch:syntactic-transition-output-eilenberg}; no additional logical
variance is introduced in this appendix.
\end{remark}

\section{Logical Synthesis}
\label{sec:logical-synthesis-formal}

The appendix can now be summarized as a single synthesis statement.  Execution
logic supplies predicates on canonical execution objects.  Behaviour
projection turns those predicates into predicates on canonical residual
behaviour.  Derivative closure and derivative interior turn behaviour
predicates into local coequational theories.  Transition--output images then
extract the associated algebraic recognizers.

\begin{theorem}[Logical synthesis theorem]
\label{thm:logical-synthesis}
Fix internal categories \(A,B\).  For every right \(B\)-action \(Y\), the
following hold.

\begin{enumerate}
  \item The canonical relative execution object is
  \[
    S^{\mathrm{can}}_Y
    =
    \mathrm{Beh}_0\times_{B_0}Y.
  \]

  \item Every Mealy morphism
  \[
    \Phi:P:A\rightsquigarrow B
  \]
  has a residual comparison map
  \[
    \beta_{P,Y}:S_{\Phi,Y}\to S^{\mathrm{can}}_Y.
  \]

  \item Primitive dynamic truth is pulled back from canonical primitive truth
  along
  \[
    \beta_{P,Y}.
  \]

  \item More generally, every formula in the residual-stable canonical
  fragment is pulled back from its canonical interpretation along the
  appropriate residual comparison map.

  \item Every local coequational theory
  \[
    V\hookrightarrow\mathrm{Beh}_0
  \]
  induces a primitive dynamic invariant
  \[
    V_Y\leq [m_{\mathrm{Beh},Y}]V_Y.
  \]

  \item Every canonical execution predicate
  \[
    \psi\hookrightarrow S^{\mathrm{can}}_Y
  \]
  has a may-coequational approximation
  \[
    \mathrm{May}_Y(\psi)
    =
    \mathrm{Der}_0(\exists_{\pi_Y}\psi),
  \]
  and a must-coequational approximation
  \[
    \mathrm{Must}_Y(\psi)
    =
    \mathrm{DerInt}(\forall_{\pi_Y}\psi).
  \]

  \item These approximations form the adjoint triple
  \[
    \mathrm{May}_Y
    \dashv
    J_Y
    \dashv
    \mathrm{Must}_Y,
    \qquad
    J_Y(V)=\pi_Y^*\,iV.
  \]

  \item Dynamic formulas, guarded response-profile formulas, finite path and
  trace formulas, and ordinary sequential projections of cubical formulas
  produce local coequational theories by applying the preceding may--must
  construction to their canonical interpretations.

  \item The generated coequational theories have syntactic
  transition--output recognizers
  \[
    \operatorname{SynTO}^{\exists/\forall}_Y(\varphi)
    =
    \operatorname{Op}
    \left(
      \mathrm{Th}^{\exists/\forall}_Y(\varphi)
    \right).
  \]

  \item The constructions are stable under downstream action maps and Mealy
  simulations, with variance determined by the corresponding downstream
  substitution maps, contravariant simulation-induced program functors, and
  the variance of transition--output images.
\end{enumerate}
\end{theorem}

\begin{proof}
\leavevmode
\begin{enumerate}
  \item \textbf{Identify canonical execution.}
  Item \(1\) is the definition of the canonical relative execution object.

  \item \textbf{Identify residual comparison.}
  Item \(2\) follows from residual semantics of Mealy morphisms.

  \item \textbf{Apply primitive pullback semantics.}
  Item \(3\) is
  Corollary~\ref{cor:primitive-canonical-pullback-semantics}.

  \item \textbf{Apply general pullback semantics.}
  Item \(4\) is Theorem~\ref{thm:canonical-pullback-semantics}.

  \item \textbf{Apply coequational invariance.}
  Item \(5\) is
  Theorem~\ref{thm:coequations-induce-primitive-invariants}.

  \item \textbf{Recall may and must approximations.}
  Item \(6\) is
  Definition~\ref{def:may-must-coequationalization}.

  \item \textbf{Use the adjoint triple.}
  Item \(7\) is
  Theorem~\ref{thm:may-must-coequational-approximation}.

  \item \textbf{Assemble the logical fragments.}
  Item \(8\) is the combination of the dynamic, guarded-profile, trace, and
  cubical-projection constructions of this appendix.

  \item \textbf{Apply transition--output images.}
  Item \(9\) is
  Definition~\ref{def:syntactic-transition-output-recognizers}.

  \item \textbf{Apply stability results.}
  Item \(10\) is Proposition~\ref{prop:downstream-substitution},
  Proposition~\ref{prop:downstream-comparison-coequational-projections}, and
  Proposition~\ref{prop:simulation-substitution-logical}, together with
  Remark~\ref{rem:recognizer-variance}.
\end{enumerate}
\end{proof}

The constructions of the appendix are summarized by the diagram
\[
\begin{tikzcd}[column sep=large,row sep=large]
\text{execution structure from Chapters 6--9}
  \arrow[r]
&
\text{logical modality}
  \arrow[r]
&
\psi\hookrightarrow X^{\mathrm{can}}
  \arrow[d, "{\exists_q,\ \forall_q}"]
\\
&
&
\text{behaviour predicate}
  \arrow[d, "{\mathrm{Der}_0,\ \mathrm{DerInt}}"]
\\
&
&
\mathbf{Coeq}_{A^\sharp,B^\sharp}
  \arrow[d, "\operatorname{Op}"]
\\
&
&
\text{syntactic transition--output recognizer}.
\end{tikzcd}
\]

For ordinary relative execution,
\[
  X^{\mathrm{can}}=S^{\mathrm{can}}_Y,
  \qquad
  q=\pi_Y,
  \qquad
  (A^\sharp,B^\sharp)=(A,B).
\]
For cubical execution, \(X^{\mathrm{can}}\) is a canonical cubical execution
object and \(q\) is one of the sequential behaviour projections supplied by
the cubical semantics.  For finite-path and feedback-trace formulas, the
path or trace span first defines a predicate transformer on
\(S^{\mathrm{can}}_Y\).  The resulting predicate on
\(S^{\mathrm{can}}_Y\) is then projected along \(\pi_Y\).  Equivalently, when
one presents the path or trace apex itself as the execution object, the
behaviour projection is the initial-state map followed by \(\pi_Y\).

Thus execution logic supplies predicates on canonical execution objects.
Quantification along the behaviour-projection map turns those predicates into
behaviour predicates.  Derivative closure and derivative interior turn the
behaviour predicates into local coequational theories.  Transition--output
images then extract the associated algebraic recognizers.  This appendix is
therefore conservative over the preceding chapters: it introduces no new
logical primitives, but records the common coequational comparison pattern
shared by the logical constructions already developed in the book.